\title{DeepSeekMath: Pushing the Limits of Mathematical Reasoning in Open Language Models}

\begin{abstract}
Mathematical reasoning remains a formidable obstacle for language models because of its highly structured and intricate nature. In this work, we present DeepSeekMath~7B, a model that extends the pre-training of DeepSeek-Coder-Base-v1.5 7B with 120 billion tokens related to mathematics, which were collected from Common Crawl, along with tokens containing natural language and code. DeepSeekMath~7B attains a remarkable 51.7\% accuracy on the MATH benchmark at the competition level, all without the assistance of external toolkits or voting strategies, thereby narrowing the gap with leading models like Gemini-Ultra and GPT-4. With self-consistency over 64 responses, DeepSeekMath~7B further boosts its MATH score to 60.9\%. The strength of DeepSeekMath~in mathematical reasoning is rooted in two principal innovations: (1) we systematically leverage the latent potential of publicly available web data through a carefully constructed data selection pipeline; (2) we propose Group Relative Policy Optimization (GRPO), an adaptation of Proximal Policy Optimization (PPO), that boosts both mathematical reasoning proficiency and the memory efficiency of PPO training.
\end{abstract}

\section{Introduction}

Large language models (LLMs) have transformed how mathematical reasoning is tackled within artificial intelligence, driving substantial progress on benchmarks for quantitative reasoning \citep{MATH} and geometric reasoning \citep{trinh2024solving}. In addition, such models are increasingly invaluable in supporting humans as they address advanced mathematical tasks \citep{tao}. Nevertheless, state-of-the-art models including GPT-4 \citep{gpt4} and Gemini-Ultra \citep{gemini} remain proprietary, and the open-source alternatives currently available show a significant performance gap in comparison.

This paper presents DeepSeekMath, a specialized language model that demonstrates substantial advancement in mathematical problem-solving compared to other open-source models and performs comparably to GPT-4 in academic benchmarks. Central to this achievement is the construction of the DeepSeekMath Corpus—a large-scale, high-quality pre-training dataset with 120B math tokens. The data is curated from Common Crawl (CC) via a fastText-based classifier \citep{joulin2016fasttext}. The process begins by training the classifier on OpenWebMath \citep{openwebmath} entries as positive samples and a broad array of miscellaneous web pages as negative samples. Using the trained classifier, additional positive examples are extracted from CC and then subjected to human annotation for further quality control. This expanded dataset is subsequently utilized to retrain and enhance the classifier. Experimental results demonstrate that the resulting dataset is of superior quality, as evidenced by DeepSeekMath-Base 7B achieving a score of 64.2\% on GSM8K \citep{gsm8k} and 36.2\% on the MATH competition benchmark \citep{MATH}, thereby surpassing Minerva 540B \citep{minerva}. Notably, the multilingual nature of DeepSeekMath Corpus also leads to improved performance on Chinese mathematics benchmarks \citep{wei2023cmath,agieval}. We suggest that our methodology in curating mathematical data provides a valuable resource for further research, while acknowledging that considerable opportunities remain for future development.

DeepSeekMath-Base is constructed by initializing with DeepSeek-Coder-Base-v1.5 7B \citep{deepseek-coder}, as empirical evidence suggests that leveraging a code-centric pretrained model yields better results than beginning with a standard large language model. In addition, we find that mathematical pretraining not only enhances the model’s proficiency in mathematics but also provides gains on general reasoning benchmarks such as MMLU \citep{mmlu} and BBH \citep{bbh}, reflecting broader improvements in cognitive abilities.

Subsequent to pre-training, we further fine-tune DeepSeekMath-Base using mathematical instruction datasets, which encompass chain-of-thought \citep{cot}, program-of-thought \citep{pot,pal}, and tool-assisted reasoning tasks \citep{tora}. This process results in the DeepSeekMath-Instruct 7B model, which outperforms other 7B parameter models and achieves performance levels rivaling instruction-tuned models with 70B parameters.

Additionally, we propose Group Relative Policy Optimization (GRPO), a modified reinforcement learning algorithm inspired by Proximal Policy Optimization (PPO) \citep{schulman2017proximal}. Unlike PPO, GRPO does not require a separate critic network; it instead infers baselines from aggregated group scores, leading to significant reductions in training cost. Using a carefully selected subset of English instruction fine-tuning data, GRPO delivers remarkable improvements over DeepSeekMath-Instruct, with notable advancements in both in-domain (GSM8K: 82.9\% $\rightarrow$ 88.2\%, MATH: 46.8\% $\rightarrow$ 51.7\%) and out-of-domain tasks (for instance, CMATH: 84.6\% $\rightarrow$ 88.8\%) during reinforcement learning. We also articulate a cohesive conceptual framework that encompasses various approaches such as Rejection Sampling Fine-Tuning (RFT) \citep{yuan2023scaling}, Direct Preference Optimization (DPO) \citep{dpo}, PPO, and GRPO, demonstrating that these techniques are unified as direct or simplified RL strategies. Through comprehensive experimentation—including comparisons of online vs. offline RL, supervision over outcomes vs. reasoning processes, and single-step vs. iterative training—we scrutinize the foundational components of this RL paradigm. Finally, we elucidate the mechanisms through which our RL framework elevates the performance of instruction-tuned models, and outline prospective research avenues for developing even more effective RL methodologies guided by this unifying framework.

\subsection{Contributions}
This work presents advancements in large-scale mathematical pre-training and delves into the evaluation and application of reinforcement learning techniques.

\textbf{Math Pre-Training at Scale}

\begin{itemize}[topsep=0pt]
    \item We demonstrate that Common Crawl's open-access web data serves as a rich source for mathematical tasks. Utilizing a rigorous data curation pipeline, we assemble the DeepSeekMath Corpus, a high-quality dataset comprising 120B tokens sourced from math-focused web content. This corpus substantially exceeds the scale of the Minerva math dataset \citep{minerva} by nearly a factor of seven, and surpasses the OpenWebMath dataset \citep{openwebmath} by roughly nine times.

    \item The DeepSeekMath-Base 7B pre-trained model achieves performance on par with Minerva 540B \citep{minerva}, suggesting that parameter count alone does not dictate success in mathematical reasoning. Our results indicate that compact models, when pre-trained with high-quality datasets, can exhibit robust mathematical capabilities.

    \item Experimental results indicate that performing code pre-training prior to math-specific training enhances a model’s proficiency in solving mathematical problems, both with and without auxiliary tools. These observations shed light on the enduring question: \textit{does code training bolster reasoning skills?} At least in the context of mathematics, our evidence supports this view.

    \item Despite widespread practice of leveraging arXiv-sourced papers for training, especially within the mathematics community, our study observes no significant gains on any mathematical benchmarks utilized in this research.
\end{itemize}

\textbf{Exploration and Analysis of Reinforcement Learning}

\begin{itemize}[topsep=0pt]
    \item This work presents Group Relative Policy Optimization (GRPO), a novel reinforcement learning algorithm characterized by both efficiency and effectiveness. GRPO eliminates the necessity for a critic model and instead leverages group-based score estimates to form a baseline, thereby drastically cutting down on the computational requirements compared to Proximal Policy Optimization (PPO).
    \item We show that GRPO yields substantial improvements for the instruction-tuned DeepSeekMath-Instruct model when utilizing only the instruction-tuning dataset. Additionally, reinforcement learning with GRPO results in measurable gains in out-of-domain generalization.
    \item Our study proposes a unified framework that accommodates various methodologies—including RFT, DPO, PPO, and GRPO—and offers a detailed experimental analysis. Investigations span online versus offline regimes, outcome-based versus process-based supervision, and comparisons between single-turn and iterative reinforcement learning, providing insight into the critical aspects underpinning this unified approach.
    \item Leveraging our unified framework, we examine the underlying factors that drive the success of reinforcement learning for language models and outline several promising avenues for enhancing the efficiency and efficacy of reinforcement learning in LLM training.
\end{itemize}

\subsection{Summary of Evaluations and Metrics}

\begin{itemize}[topsep=0pt]
    \item \textbf{English and Chinese Mathematical Reasoning}: We carry out rigorous evaluations of our models on a variety of English and Chinese benchmarks that encompass mathematics problems ranging from elementary to undergraduate levels. The English datasets assessed include GSM8K \citep{gsm8k}, MATH \citep{MATH}, SAT \citep{llemma}, OCW Courses \citep{minerva}, and MMLU-STEM \citep{mmlu}; for Chinese, we use MGSM-zh \citep{mgsm}, CMATH \citep{wei2023cmath}, Gaokao-MathCloze \citep{agieval}, and Gaokao-MathQA \citep{agieval}. The evaluation focuses both on the models’ ability to compose independent textual solutions and on their problem-solving capabilities via Python code.

    In the English context, DeepSeekMath-Base demonstrates performance that rivals the proprietary Minerva 540B \citep{minerva} and outperforms all open-source base models—such as Mistral 7B \citep{mistral} and Llemma-34B \citep{llemma}—irrespective of the presence of mathematical pre-training, often with notable margins. For the Chinese benchmarks, DeepSeekMath-Base achieves even stronger results. We attribute this to our choice of incorporating high-quality multilingual data, diverging from earlier approaches \citep{minerva, llemma} which limited themselves to English-only sources. After applying instruction tuning and reinforcement learning, both DeepSeekMath-Instruct and DeepSeekMath-RL deliver state-of-the-art outcomes, becoming the first open-source models to achieve over 50\% accuracy on the highly challenging, competition-level MATH dataset.
\end{itemize}

\item \textbf{Formal Mathematics}: The performance of DeepSeekMath-Base is assessed on the informal-to-formal theorem proving benchmark introduced in \citep{dsp_proof}, specifically utilizing the miniF2F dataset \citep{minif2f} with Isabelle \citep{isabelle} serving as the proof assistant. Results indicate that DeepSeekMath-Base achieves notable few-shot capabilities in the task of autoformalization.

\item \textbf{Natural Language Understanding, Reasoning, and Code}: To achieve a holistic evaluation of the models' capabilities across understanding, reasoning, and programming, DeepSeekMath-Base is benchmarked on the Massive Multitask Language Understanding (MMLU) dataset \citep{mmlu}, which features 57 multiple-choice problems spanning a wide array of domains; on BIG-Bench Hard (BBH) \citep{bbh}, composed of 23 particularly challenging tasks emphasizing multi-step reasoning; as well as on HumanEval \citep{codex} and MBPP \citep{mbpp}, both established benchmarks for evaluating code generation models. Empirical results suggest that pre-training on mathematical corpora enhances both linguistic comprehension and logical reasoning abilities.

\section{Math Pre-Training}

\subsection{Data Collection and Decontamination} This subsection describes the construction pipeline for the DeepSeekMath Corpus using Common Crawl as a primary source. Illustrated in Figure \ref{fig:data_collect}, our iterative process details how a large-scale mathematical dataset can be extracted methodically from Common Crawl, initiating with a carefully selected seed set consisting of high-quality math-related material. Notably, this framework can be extended to other specialized domains, such as computer programming.

Initially, we adopt OpenWebMath \citep{openwebmath}, a curated collection of high-quality mathematics web data, as the foundational seed dataset. This collection is then leveraged to train a fastText classifier \citep{joulin2016fasttext} with the objective of identifying additional web pages similar in content to OpenWebMath. Concretely, 500,000 samples from the seed set are labeled as positive instances, while another 500,000 randomly selected web pages from Common Crawl act as negatives. The classifier is implemented using a public library\footnote{\url{https://fasttext.cc}}, with settings as follows: 256-dimensional vectors, learning rate of 0.1, a maximum word n-gram size of 3, a minimum word frequency threshold of 3, and 3 epochs for training. To narrow down the initial scale of Common Crawl, both URL-based deduplication and near-deduplication are employed, producing a filtered set of 40B unique HTML pages. The fastText classifier is then applied to these pages to extract mathematically relevant documents. To ensure dataset quality, collected pages are scored via the classifier, and only those with the highest rankings are retained. The retention threshold is informed by pre-training trials across token subsets of 40B, 80B, 120B, and 160B; for the first iteration, we maintain the top 40B tokens.

Following the initial round of data collection, a substantial number of mathematical web pages are not yet included, largely due to the limited diversity within the positive samples used to train the fastText model. To address this shortcoming, we expand the seed corpus by identifying additional math-focused web resources, thus refining the model’s effectiveness. The process begins by segmenting all Common Crawl data into non-overlapping domains, where each domain corresponds to web pages sharing the same base URL. We then compute, for every domain, the proportion of its pages retrieved during the first collection pass. Domains with more than 10\% of their pages collected are flagged as math-oriented (e.g., \url{mathoverflow.net}).

We next conduct manual annotation to identify URLs in these selected domains that host mathematical content (e.g., \url{mathoverflow.net/questions}). Those uncollected pages associated with these annotated URLs are subsequently incorporated into the seed corpus. This methodology enriches the pool of positive training examples, facilitating the training of a more effective fastText model that can recover additional mathematical web pages in the next iteration. Upon completing four collection rounds, the resulting corpus comprises 35.5 million mathematical web pages, totaling 120 billion tokens. By the fourth round, nearly 98\% of pages have already been gathered by the third round, prompting us to discontinue further collection.

To mitigate the risk of benchmark data leakage, we adopt the filtering procedures described by \cite{deepseek-coder}, excluding web pages that include questions or answers from widely used English math benchmarks like GSM8K \citep{gsm8k} and MATH \citep{MATH}, as well as Chinese benchmarks such as CMATH \citep{wei2023cmath} and AGIEval \citep{agieval}. Specifically, any text fragment containing a 10-gram sequence that is an exact match for a substring in the benchmark datasets is removed from the training set. For benchmark segments that are shorter than 10 grams but at least 3 grams long, we also use exact string matching to exclude any contaminated material.

\subsection{Validating the Quality of the DeepSeekMath~Corpus}

To assess the quality of the DeepSeekMath~Corpus, we conduct pre-training experiments, benchmarking it against other prominent recently introduced math corpora:

\begin{itemize}[topsep=0pt]
    \item \textbf{MathPile} \citep{mathpile}: an 8.9B-token composite dataset assembled from sources such as textbooks, Wikipedia, ProofWiki, CommonCrawl, StackExchange, and arXiv, where arXiv constitutes more than 85\% of the content;
    \item \textbf{OpenWebMath} \citep{openwebmath}: a subset of CommonCrawl comprising 13.6B tokens, filtered for math content;
    \item \textbf{Proof-Pile-2} \citep{llemma}: a corpus constructed from OpenWebMath, AlgebraicStack (offering 10.3B tokens of math code), and arXiv articles (28.0B tokens). For experiments with Proof-Pile-2, we adhere to the data ratio of arXiv:Web:Code as 2:4:1 as detailed in \cite{llemma}.
\end{itemize}

\subsubsection{Training Setting}
\label{sec:quality-policy}
We perform mathematical domain adaptation on a general-purpose language model with 1.3B parameters, structurally aligned with the DeepSeek LLMs \citep{deepseek-llm}, henceforth referred to as DeepSeek-LLM 1.3B. Each mathematical dataset is used to fine-tune a separate model for a total of 150B training tokens. All training runs utilize the resource-efficient and streamlined HAI-LLM \citep{haillm} framework. In keeping with the established DeepSeek LLM training procedures, we adopt the AdamW optimizer \citep{adamW} configured with $\beta_1=0.9$, $\beta_2=0.95$, and $\mathrm{weight\_decay}=0.1$. The learning rate follows a multi-phase decay schedule: it initially ramps up to its maximum value over 2,000 warmup steps, then declines to 31.6\% of its peak after 80\% of the training duration, and further drops to 10.0\% of its initial maximum by the completion of 90\% of training. The peak learning rate is set at 5.3e-4. We employ a batch size of 4 million tokens and a context window of 4,000 tokens.

\subsubsection{Evaluation Results}

\textbf{The DeepSeekMath~Corpus exhibits superior quality, robust multilingual coverage, and is the most extensive dataset in its category.}

\begin{itemize}[topsep=0pt]
    \item \textbf{High-quality}:
    We assess downstream capabilities across eight mathematical benchmark tasks using few-shot chain-of-thought prompting \cite{cot}. According to Table \ref{tab:corpora_comparison}, models fine-tuned on the DeepSeekMath~Corpus consistently outperform counterparts trained on alternative datasets. As visualized in Figure \ref{fig:corpora_comparisons}, models trained with DeepSeekMath~Corpus surpass the performance of those trained on Proof-Pile-2 after 50B tokens (which corresponds to one full epoch on Proof-Pile-2), illustrating the superior average data quality in DeepSeekMath~Corpus.
    \item \textbf{Multilingual}:
    The DeepSeekMath~Corpus integrates mathematical text from a variety of languages, with English and Chinese comprising the majority. Table \ref{tab:corpora_comparison} demonstrates that leveraging DeepSeekMath~Corpus for training enhances mathematical reasoning capabilities in both English and Chinese. In contrast, legacy mathematical corpora—largely dominated by English—either yield only minimal gains or even negatively affect mathematical reasoning performance in Chinese.
    \item \textbf{Large-scale}:
    Relative to prior mathematical corpora, the DeepSeekMath~Corpus is substantially larger in scope. As seen in Figure \ref{fig:corpora_comparisons}, when DeepSeek-LLM 1.3B is trained on DeepSeekMath~Corpus, it exhibits a sharper and more prolonged performance improvement trajectory. Conversely, baseline datasets, being considerably smaller, are cycled through multiple times during training, which quickly leads to stagnation in model performance.
\end{itemize}

\subsection{Training and Evaluating DeepSeekMath-Base 7B}

This section presents DeepSeekMath-Base 7B, a foundational model designed to demonstrate advanced mathematical reasoning. We initialized this model with DeepSeek-Coder-Base-v1.5 7B \citep{deepseek-coder} and trained it on a dataset consisting of 500B tokens. The data distribution includes 56\% from the DeepSeekMath Corpus, 4\% from AlgebraicStack, 10\% sourced from arXiv, 20\% consisting of Github code, and the final 10\% drawn from natural language data in both English and Chinese, collected from Common Crawl. Our training configuration largely follows the procedure described in Section \ref{sec:quality-policy}, with the distinction that the learning rate peaks at 4.2e-4 and a batch size of 10M tokens is used.

A thorough evaluation of DeepSeekMath-Base 7B is conducted to measure its mathematical reasoning skills. Specifically, we examine its capacity to generate independent mathematical solutions, its proficiency in solving math tasks with the assistance of computational tools, and its formal theorem proving abilities. In addition, we broaden the assessment to include general capabilities such as natural language understanding, logical reasoning, and programming.

\paragraph{Mathematical Problem Solving with Step-by-Step Reasoning}

We assess DeepSeekMath-Base’s aptitude for mathematical problem solving through few-shot chain-of-thought prompting \citep{cot} on eight benchmarks in both English and Chinese. These benchmarks address a range of mathematical challenges, including quantitative reasoning tasks (e.g., GSM8K \citep{gsm8k}, MATH \citep{MATH}, and CMATH \citep{wei2023cmath}) as well as multiple-choice questions (e.g., MMLU-STEM \citep{mmlu}, Gaokao-MathQA \citep{agieval}), spanning mathematical domains from basic arithmetic to university-level topics.

As presented in Table \ref{tab:base_math_cot}, DeepSeekMath-Base 7B establishes state-of-the-art results across all eight benchmark tasks when compared with other open-source base models, such as the well-known Mistral 7B \citep{mistral} and the more recent Llemma 34B \citep{llemma}, which benefited from math-specific training on Proof-Pile-2 \citep{llemma}. In particular, DeepSeekMath-Base achieves an absolute improvement exceeding 10\% over all existing open-source base models on the challenging MATH competition dataset, and even surpasses the closed-source Minerva 540B \citep{minerva}—a model based on PaLM \citep{palm} that is 77 times larger and specialized through additional mathematical corpora.

\paragraph{Mathematical Problem Solving with Tool Use} To assess program-based mathematical reasoning, we conducted evaluations on GSM8K and MATH datasets, employing few-shot program-of-thought prompting \citep{pot,pal}. Each model receives a prompt to tackle a problem by writing a Python program, making use of libraries such as \textit{math} and \textit{sympy} for more complex calculations. The program’s output is considered as the solution. Results in Table \ref{tab:base_math_pot_proof} demonstrate that DeepSeekMath-Base 7B achieves superior performance to the previously leading Llemma 34B.

\paragraph{Formal Mathematics}

Automating formal proof generation is crucial for both ensuring correctness and boosting the productivity of mathematical reasoning, drawing increasing research focus. DeepSeekMath-Base 7B was assessed on the informal-to-formal proof generation task described in \citep{dsp_proof}, where the model must construct a formal proof from a provided informal statement, its formal restatement, and an informal justification. Using the miniF2F \citep{minif2f} benchmark—which targets Olympiad-level formal mathematics—each instance required the generation of a formal Isabelle proof using few-shot prompts. Following procedures from \cite{dsp_proof}, the model initially generates a sketch of the proof, and Sledgehammer \citep{sledgehammer} is then used to complete any remaining proof steps. As indicated in Table \ref{tab:base_math_pot_proof}, DeepSeekMath-Base 7B exhibits strong results in proof autoformalization.

\paragraph{Natural Language Understanding, Reasoning, and Code}

Model performance in natural language understanding, complex reasoning, and coding was measured on MMLU \citep{mmlu}, BBH \citep{bbh}, HumanEval \citep{codex}, and MBPP \citep{mbpp}. According to Table \ref{tab:understanding_reasoning_code}, DeepSeekMath-Base 7B achieves marked improvements on both MMLU and BBH compared to its predecessor DeepSeek-Coder-Base-v1.5 \citep{deepseek-coder}, underscoring the advantages of math-oriented training for language understanding and logical reasoning. Furthermore, the integration of code tokens during continual pretraining allows DeepSeekMath-Base 7B to retain competitive performance on the coding tasks relative to DeepSeek-Coder-Base-v1.5. Overall, DeepSeekMath-Base 7B delivers substantially stronger results on the three reasoning and coding benchmarks when compared to the general-purpose Mistral 7B \citep{mistral}.

\section{Supervised Fine-Tuning}

\subsection{SFT Data Curation}

We compile an instruction-tuning dataset for mathematics that spans both English and Chinese languages, encompassing a range of mathematical domains and varying difficulty. Each problem is paired with solutions represented as chain-of-thought (CoT) \citep{cot}, program-of-thought (PoT) \citep{pot,pal}, or tool-integrated reasoning \citep{tora}. The assembled dataset comprises a total of 776K training samples.

\begin{itemize}[topsep=0pt]
    \item \textbf{English mathematical datasets}: We provide tool-integrated solutions for GSM8K and MATH datasets, and include a selected portion of MathInstruct \citep{MathInstruct} together with the training split of Lila-OOD \citep{lila}, in which CoT or PoT solutions are employed. The English dataset incorporates a wide array of mathematical disciplines such as algebra, probability, number theory, calculus, and geometry.
    \item \textbf{Chinese mathematical datasets}: For the Chinese collection, we gather K-12 mathematical problems distributed across 76 sub-categories (e.g., linear equations), with corresponding solutions annotated using both CoT and tool-integrated approaches.
\end{itemize}

\subsection{Training and Evaluating DeepSeekMath-Instruct 7B}

Here, we describe DeepSeekMath-Instruct 7B, which is refined via instruction tuning built on the DeepSeekMath-Base model. Training data samples are concatenated at random, up to a context window of 4K tokens. The training regime consists of 500 optimization steps, using a batch size of 256 and a fixed learning rate of 5e-5.

We assess mathematical proficiency of models in both tool-augmented and tool-free settings, utilizing four quantitative reasoning benchmarks in English and Chinese. Comparative evaluation is conducted against top-tier contemporary models:

\begin{itemize}[topsep=0pt]
    \item \textbf{Closed-source models}: These encompass (1) models from the GPT suite, notably GPT-4 \citep{gpt4} and GPT-4 Code Interpreter~\footnote{\url{https://openai.com/blog/chatgpt-plugins\#\#code-interpreter}}, (2) Gemini Ultra and Pro \citep{gemini}, (3) Inflection-2 \citep{inflection-2}, (4) Grok-1~\footnote{\url{https://x.ai/model-card}}, as well as recent Chinese offerings such as (5) Baichuan-3~\footnote{\url{https://www.baichuan-ai.com}}, and (6) the most recent model from the GLM series, GLM-4~\footnote{\url{https://open.bigmodel.cn/dev/api\#glm-4}} \citep{glm}.
\end{itemize}

These models serve general-purpose applications and have typically been subjected to comprehensive alignment processes.
    \item \textbf{Open-source models} in this study encompass: general-purpose systems such as (1) DeepSeek-LLM-Chat 67B \citep{deepseek-llm}, (2) Qwen 72B \citep{qwen}, (3) SeaLLM-v2 7B \citep{seallm}, and (4) ChatGLM3 6B \citep{chatglm3}. In addition, several models with targeted improvements for mathematical reasoning are included: (5) InternLM2-Math 20B~\footnote{\url{https://github.com/InternLM/InternLM-Math}}, which extends InternLM2 with both mathematics-oriented training and subsequent instruction tuning; (6) Math-Shepherd-Mistral 7B, utilizing PPO optimization \citep{schulman2017proximal} applied to Mistral 7B \citep{mistral} in conjunction with a process-supervised reward framework; (7) the WizardMath series \citep{wizardmath}, which augments mathematical abilities in Mistral 7B and Llama-2 70B \citep{llama2} through the evolve-instruct strategy (a variant of instruction tuning employing machine-generated tasks) and PPO, largely drawing training data from GSM8K and MATH; (8) MetaMath 70B \citep{metamath}, which is based on Llama-2 70B with fine-tuning on extended GSM8K and MATH datasets; (9) ToRA 34B \cite{tora}, developed from CodeLlama 34B and fine-tuned for integrated mathematical tool reasoning; and (10) MAmmoTH 70B \citep{MathInstruct}, where Llama-2 70B is instruction-tuned using MathInstruct.

Referencing Table \ref{tab:sft_rl_math}, in the evaluation scenario prohibiting the use of external tools, DeepSeekMath-Instruct 7B achieves notable step-by-step reasoning capabilities. Particularly on the competition-standard MATH dataset, this model exceeds the accuracy of all surveyed open-source alternatives as well as most proprietary offerings (such as Inflection-2 and Gemini Pro) by margins exceeding 9\% in absolute terms. This advantage holds even in comparison to considerably larger models (e.g., Qwen 72B) and models that have undergone extensive reinforcement learning on mathematical tasks (for example, WizardMath-v1.1 7B). While DeepSeekMath-Instruct matches the performance of leading proprietary Chinese models like GLM-4 and Baichuan-3 on MATH, it does not yet surpass GPT-4 or Gemini Ultra.

In an alternative evaluation configuration permitting both natural language reasoning and code-based tool use, DeepSeekMath-Instruct 7B attains nearly 60\% accuracy on MATH, outperforming all other open-source models in the literature. On additional benchmarks, its results are on par with those of DeepSeek-LLM-Chat 67B, which was previously considered state-of-the-art and operates at roughly tenfold the scale.

\section{Reinforcement Learning}

\subsection{Group Relative Policy Optimization} Reinforcement learning (RL) has shown considerable efficacy in enhancing the mathematical reasoning skills of LLMs after they undergo Supervised Fine-Tuning (SFT) \citep{wang2023math,wizardmath}. Here, we present a novel and efficient RL technique, Group Relative Policy Optimization (GRPO).

\subsubsection{From PPO to GRPO} Proximal Policy Optimization (PPO) \citep{schulman2017proximal} is a well-established actor-critic RL technique frequently utilized for RL fine-tuning of LLMs \citep{ouyang2022training}. The method seeks to optimize LLMs by maximizing the surrogate objective given by:

\begin{equation}
\footnotesize
    \mathcal{J}_{PPO}(\theta) = \mathbb{E}{[q \sim P(Q), o \sim \pi_{\theta_{old}}(O|q)]} \frac{1}{|o|} \sum_{t=1}^{|o|} \min \left[ \frac{\pi_\theta(o_{t} | q, o_{<t})}{\pi_{\theta_{old}}(o_{t} | q, o_{<t})} A_{t}, \text{clip} \left( \frac{\pi_\theta(o_{t} | q, o_{<t})}{\pi_{\theta_{old}}(o_{t} | q, o_{<t})},

In this context, $\pi_{\theta}$ and $\pi_{\theta_{old}}$ denote the present and prior policy models, respectively, while $q$ and $o$ represent questions and outputs sampled either from the question corpus or from $\pi_{\theta_{old}}$. The parameter $\varepsilon$ is a hyper-parameter introduced by PPO to control clipping for improved training robustness. The advantage term $A_t$ is evaluated using Generalized Advantage Estimation (GAE) \citep{gae}, computed from the reward sequence $\{r_{\ge t}\}$ and a value function $V_{\psi}$ that is learned during training. Consequently, PPO requires not only the policy model but also a value function model to be trained simultaneously. To guard against excessive reward model optimization, a common approach appends a per-token KL penalty—computed between the policy and a reference model—to each token’s reward \citep{ouyang2022training}, leading to the following formulation:

\begin{equation}
    r_{t} = r_\varphi(q, o_{\le t}) - \beta \log\frac{\pi_{\theta}(o_{t}|q, o_{<t})}{\pi_{ref}(o_{t}|q, o_{<t})},
\label{eq:PPO-reward}
\end{equation}

Here, $r_\varphi$ stands for the reward model, $\pi_{ref}$ refers to the reference policy, typically chosen as the initial SFT model, and $\beta$ scales the KL penalty term.

The use of a separate value function in PPO generally implies maintaining an auxiliary network of comparable size to the policy model, which substantially increases resource demands for both computation and memory. Moreover, the value function is mainly intended to serve as a baseline for variance reduction in the computation of advantages during RL updates. In LLM applications, the reward model customarily assigns a score only to the final output token, which can hinder effective value function training at the token level. To overcome these challenges, we introduce Group Relative Policy Optimization (GRPO), depicted in Figure \ref{fig:grpo}. Unlike PPO, GRPO does not require an explicit value function; instead, it computes a baseline using the mean reward obtained from multiple candidate outputs sampled for the same question. Concretely, for each question $q$, we generate a collection of outputs $\{o_1, o_2, \cdots, o_G\}$ by sampling from the previous policy $\pi_{\theta_{old}}$, and subsequently update the policy by maximizing the following objective:

\begin{equation}
\footnotesize
\begin{split}
    \mathcal{J}_{GRPO}(\theta) &= \mathbb{E}{[q \sim P(Q), \{o_i\}_{i=1}^G \sim \pi_{\theta_{old}}(O|q)]}  \\
    & \frac{1}{G}\sum_{i=1}^G\frac{1}{|o_i|} \sum_{t=1}^{|o_i|} \left\{ \min \left[ \frac{\pi_\theta(o_{i,t} | q, o_{i,<t})}{\pi_{\theta_{old}}(o_{i,t} | q, o_{i,<t})} \hat{A}_{i,t}, \text{clip} \left( \frac{\pi_\theta(o_{i,t} | q, o_{i,<t})}{\pi_{\theta_{old}}(o_{i,t} | q, o_{i,<t})}, 1 - \varepsilon, 1 + \varepsilon \right)  \hat{A}_{i,t} \right] - \beta \mathbb{D}_{KL}\left[\pi_{\theta} || \pi_{ref}\right]\right\} ,
\end{split}
\label{eq:GRPO-obj}
\end{equation}

In this formulation, $\varepsilon$ and $\beta$ remain as hyper-parameters, while $\hat{A}_{i,t}$ captures the group-relative advantage for each output, computed using only the rewards from within the same group—a detail elaborated upon in later subsections. This group-based advantage estimation closely aligns with the nature of reward models, which are typically designed to assess outputs by direct comparison for the same prompt. Furthermore, instead of incorporating the KL term into the reward itself, GRPO introduces regularization by applying the KL divergence between the learned policy and the reference policy directly to the loss, thus streamlining the calculation of $\hat{A}_{i,t}$. Unlike the per-token KL penalty applied in (\ref{eq:PPO-reward}), we employ an unbiased estimator for KL divergence, as described in

\begin{algorithm}[t]
  \small
  \caption{Iterative Group Relative Policy Optimization}
  \textbf{Input} initial policy model $\pi_{\theta_{\text{init}}}$; reward models $r_\varphi$; task prompts $\mathcal{D}$; 
  hyperparameters $\varepsilon$, $\beta$, $\mu$
  \begin{algorithmic}[1]
    \State Initialize policy model: $\pi_\theta \leftarrow \pi_{\theta_{\text{init}}}$
    \For{iteration = 1, \dots, I}
      \State Set reference model $\pi_{ref} \leftarrow \pi_\theta$
      \For{step = 1, \dots, M}
        \State Select batch $\mathcal{D}_b$ from dataset $\mathcal{D}$
        \State Update the previous policy: $\pi_{\theta_{old}} \leftarrow \pi_{\theta}$
        \State For each question $q \in \mathcal{D}_b$, sample $G$ responses $\{o_i\}_{i=1}^G \sim \pi_{\theta_{old}} (\cdot \mid q)$
        \State Assign reward values $\{r_i\}_{i=1}^{G}$ to each $o_i$ using $r_{\varphi}$
        \State Determine $\hat{A}_{i,t}$ at token $t$ in $o_i$ via group-based relative advantage estimation
        \For{GRPO iteration = 1, \dots, $\mu$}
          \State Adjust policy model $\pi_{\theta}$ to maximize the GRPO objective (see Equation \ref{eq:GC-GRPO})
        \EndFor
      \EndFor
      \State Periodically update $r_\varphi$ through ongoing training using experience replay.
    \EndFor
  \end{algorithmic}
  \textbf{Output} $\pi_\theta$
  \label{alg:iter-grpo}
\end{algorithm}

\subsubsection{Outcome Supervision RL with GRPO} More specifically, for a given question $q$, one draws a set of responses $\{o_1, o_2, \cdots, o_G\}$ from the prior policy $\pi_{\theta_{old}}$. These outputs are then evaluated by a reward model, resulting in a set of scores $\mathbf{r} = \{r_1, r_2, \cdots, r_G\}$. To ensure comparability, the raw reward scores are normalized by subtracting the mean and dividing by the standard deviation across the group. Under outcome supervision, the normalized reward for the entire output $o_i$ is uniformly assigned as the advantage $\hat{A}_{i,t}$ for each token $t$ within that output; specifically, $\hat{A}_{i, t} = \widetilde{r}_i = \frac{r_i- {\rm mean}(\mathbf{r})}{{\rm std}(\mathbf{r})}$. The optimization of the policy is performed to maximize the objective presented in equation (\ref{eq:GRPO-obj}).

\subsubsection{Process Supervision RL with GRPO} Relying solely on outcome-level rewards may provide limited and delayed feedback, which can hinder learning effectiveness in complex mathematical reasoning tasks. Building on the approach from \cite{wang2023math}, we consider process supervision, in which rewards are provided after each intermediate reasoning step. Given a question $q$ and a set of sampled responses $\{o_1, o_2, \cdots, o_G\}$, a dedicated process reward model evaluates each step within the outputs, generating rewards $\mathbf{R} = \{ \{r_1^{index(1)},\cdots,r_1^{index(K_1)}\}, \cdots,  \{r_G^{index(1)},\cdots,r_G^{index(K_G)}\} \}$, where $index(j)$ denotes the end-of-step token for the $j$-th step and $K_i$ counts the total steps in output $i$. These step-level rewards are also standardized via subtraction of the mean and division by the standard deviation: $\widetilde{r}_i^{index(j)} = \frac{r_i^{index(j)} - {\rm mean(\mathbf{R})}}{{\rm std(\mathbf{R})}}$. The advantage attributed to token $t$ in output $o_i$ is then calculated as the cumulative sum of all subsequent normalized step rewards, i.e., $\hat{A}_{i, t} = \sum_{index(j) \ge t} \widetilde{r}_i^{index(j

\subsubsection{Iterative RL with GRPO} As the reinforcement learning progresses, the existing reward model may no longer adequately guide the evolving policy model. To address this, we investigate iterative RL utilizing GRPO. Algorithm \ref{alg:iter-grpo} outlines this process: in each iteration, new training data for the reward model are produced using samples generated by the current policy model. The reward model is updated by leveraging a replay mechanism, which mixes in 10\% of previously collected data. Subsequently, the reference model is aligned with the updated policy model, which is further trained using the revised reward model.

\subsection{Training and Evaluating DeepSeekMath-RL}

Our RL experiments are conducted on DeepSeekMath-Instruct 7B. The RL training corpus comprises chain-of-thought-style questions drawn from GSM8K and MATH subsets within the SFT data, totaling approximately 144K samples. To isolate the effect of RL on datasets that lack RL-phase exposure, we omit other SFT questions. The procedure for assembling reward model training data follows the methodology from \citep{wang2023math}. The reward model is initialized from DeepSeekMath-Base 7B and trained with a learning rate of 2e-5. When applying GRPO, the policy model uses a learning rate of 1e-6, and a KL penalty coefficient of 0.04. For every question, we generate $64$ completions. The sequence length is capped at 1024, with a batch size of 1024. The policy model receives just one update after each sampling round. For assessment, DeepSeekMath-RL 7B is evaluated on the same benchmarks as DeepSeekMath-Instruct 7B. Tasks from GSM8K and MATH that involve chain-of-thought reasoning are categorized as in-domain, whereas the remaining benchmarks are considered out-of-domain.

Table \ref{tab:sft_rl_math} presents a comparative analysis of open-source and closed-source models, evaluating their performance on both chain-of-thought and tool-augmented reasoning tasks across English and Chinese datasets. Our observations are as follows: 1) DeepSeekMath-RL 7B achieves accuracy scores of 88.2\% on GSM8K and 51.7\% on MATH through chain-of-thought-based reasoning, outperforming all open-source counterparts in the 7B–70B parameter range and exceeding the capabilities of most closed-source alternatives. 2) Notably, DeepSeekMath-RL 7B is exclusively trained with chain-of-thought instruction data from GSM8K and MATH, building upon DeepSeekMath-Instruct 7B as its base. In spite of this limited training corpus, the model consistently exceeds DeepSeekMath-Instruct 7B in every evaluated category, highlighting the substantial gains introduced by reinforcement learning.

\section{Discussion}
This section provides an overview of insights derived from both pre-training and reinforcement learning experiments.

\subsection{Lessons Learnt in Pre-Training}

Here, we elaborate on the knowledge gained throughout the pre-training phase. Unless explicitly mentioned otherwise, all procedures conform to the methodologies discussed in Section \ref{sec:quality-policy}. Furthermore, throughout this section, any reference to the DeepSeekMath Corpus pertains to an 89B-token collection, compiled during the project’s second round of data acquisition.

\subsubsection{Code Training Benefits Mathematical Reasoning}

A widely-held, though not thoroughly validated, belief is that incorporating code into model training enhances reasoning ability. Our findings provide partial support for this hypothesis within the mathematical context: integrating code training bolsters models’ mathematical reasoning skills, whether or not external tools are leveraged.

To investigate the effect of code training on mathematical reasoning, we implemented and assessed both a two-stage training paradigm and a single-stage training protocol:

\textbf{Two-Stage Training}

\begin{itemize}[topsep=0pt]
    \item \textbf{Code Training for 400B Tokens $\rightarrow$ Math Training for 150B Tokens}: DeepSeek-LLM 1.3B undergoes training on 400B tokens of code prior to further training with 150B tokens of mathematical data;
    \item \textbf{General Training for 400B Tokens $\rightarrow$ Math Training for 150B Tokens}: As a comparison, in another configuration, we substitute the code tokens in the initial phase with general tokens (extracted from a large-scale general corpus assembled by DeepSeek-AI). This experiment aims to assess whether pretraining on code tokens provides an edge over general tokens in advancing mathematical reasoning skills.
\end{itemize}

\textbf{One-Stage Training}

\begin{itemize}[topsep=0pt]
    \item \textbf{Math Training for 150B Tokens}: DeepSeek-LLM 1.3B is trained exclusively on 150B math tokens;
    \item \textbf{Training on a mixture of 400B Code Tokens and 150B Math Tokens}: Since mathematical training following code pretraining tends to impair coding ability, we examine the effects of blending code and math tokens within a single-stage training regimen to determine if this approach simultaneously enhances mathematical reasoning and helps prevent catastrophic forgetting.
\end{itemize}

\paragraph{Results}
The downstream outcomes for these various training paradigms are presented in Table \ref{tab:code-math} and Table \ref{tab:code-math-general-eval}.

Incorporating code pretraining yields measurable improvements in program-aided mathematical reasoning across both staged and mixed-training protocols. According to Table \ref{tab:code-math}, two-stage pretraining with code significantly augments performance on GSM8K and MATH datasets when employing Python-based approaches. Introducing a subsequent math-specific training phase brings about further performance gains. Notably, within the one-stage mixed training setting, jointly processing code and math tokens helps reduce catastrophic forgetting typical of two-stage approaches, while fostering synergy in both coding (Table \ref{tab:code-math-general-eval}) and mathematical problem solving (Table \ref{tab:code-math}).

Beyond tool-assisted contexts, code pretraining still provides advantages for mathematical reasoning. In the sequential (two-stage) setup, starting with code-focused learning delivers appreciable boosts and also enhances the effectiveness of later math-centric training, culminating in peak results. In contrast, the simultaneous (one-stage) presentation of code and math tokens appears to detract from tool-free mathematical reasoning abilities. We hypothesize that DeepSeek-LLM 1.3B's restricted model capacity constrains its ability to efficiently absorb both programming and mathematical content concurrently.

\subsubsection{ArXiv Papers Appear Unhelpful for Enhancing Mathematical Reasoning}

Math pre-training datasets frequently incorporate arXiv papers as a key resource \citep{minerva,gpt-f,llemma,mathpile}. Nonetheless, their actual influence on the development of mathematical reasoning abilities has not been scrutinized in depth. Contrary to what one might expect, our empirical findings indicate that arXiv papers do not substantially contribute to improvements in mathematical reasoning. We conducted experiments across different model scales—including DeepSeek-LLM 1.3B and DeepSeek-Coder-Base-v1.5 7B \citep{deepseek-coder}—leveraging various processed versions of the arXiv corpus:

\begin{itemize}[topsep=0pt]
    \item \textbf{MathPile} \citep{mathpile}: an 8.9B-token dataset curated using specific cleaning and heuristic filtering strategies, with over 85\% of the content sourced from scientific arXiv papers;
    \item \textbf{ArXiv-RedPajama} \citep{redpajama}: a 28.0B-token collection comprised of raw arXiv LaTeX documents from which preambles, comments, macros, and references have been excised.
\end{itemize}

In the study, DeepSeek-LLM 1.3B was trained on each arXiv dataset for 150B tokens and DeepSeek-Coder-Base-v1.5 7B for 40B tokens. Results consistently reveal that exposure to arXiv-only material does not yield discernible progress, and may even result in reduced performance, on a diverse set of mathematical benchmarks ranging in complexity. These assessments encompass datasets for quantitative reasoning such as GSM8K and MATH (see Table \ref{tab:arxiv-cot}), multiple-choice tasks like MMLU-STEM (see Table \ref{tab:arxiv-cot}), and formal proof challenges like miniF2F (see Table \ref{tab:arxiv_atp}).

Yet, these observations are bounded by certain caveats and must be interpreted cautiously. Specifically, we have not yet examined:

\begin{itemize}[topsep=0pt]
    \item The role of arXiv data in niche mathematical tasks outside our evaluation scope, for instance, translating formal statements or proofs into their informal counterparts (the informalization problem);
    \item The interaction effects when arXiv material is mixed with other training data;
    \item Whether benefits might emerge with substantially larger model architectures.
\end{itemize}

Accordingly, more in-depth research is warranted, and we defer these questions to subsequent investigations.

\subsection{Insights from Reinforcement Learning}
\subsubsection{Towards a Unified Paradigm} This section aims to synthesize a generalizable framework for understanding various training algorithms, including SFT, RFT, DPO, PPO, and GRPO, while also presenting experimental studies probing components of this unified structure. Broadly, the parameter gradient $\theta$ for these methods can be represented as follows:

\begin{equation}
    \nabla_{\theta}\mathcal{J}_{\textcolor{red}{\mathcal{A}}}(\theta) = \mathbb{E}[\underbrace{(q,o) \sim \textcolor{red}{\mathcal{D}}}_{Data \ Source}]\left( \frac{1}{|o|} \sum_{t=1}^{|o|}  \underbrace{GC_{{\mathcal{A}}}(q, o, t, \textcolor{red}{\pi_{{rf}}})}_{Gradient \ Coefficient}  \nabla_{\theta}\log \pi_{\theta}(o_t | q, o_{<t})\right).
\label{eq:objective}
\end{equation}

This formulation is structured around three fundamental elements: 1) the \textit{Data Source} $\mathcal{D}$, which specifies the origin of training examples; 2) the \textit{Reward Function} $\pi_{{rf}}$, serving as the mechanism that produces the reward signal for training; and 3) the \textit{Algorithm} $\mathcal{A}$, which combines the training samples and the reward signal to generate the gradient coefficient $GC$, thus dictating the extent of positive or negative updates for the parameters. Within this unified analytical framework, we examine several prototypical approaches:

\begin{itemize}
    \item  \textbf{Supervised Fine-tuning (SFT)}: SFT further adapts a pretrained model using human-curated SFT datasets.
    \item \textbf{Rejection Sampling Fine-tuning (RFT)}: RFT extends SFT by re-training the model on a subset of generated outputs, selected for their correctness, with both samples and questions derived from the SFT model.
    \item \textbf{Direct Preference Optimization (DPO)}: DPO builds on SFT by employing a pair-wise preference loss for training on model-augmented responses to enhance alignment with human preferences.
    \item \textbf{Online Rejection Sampling Fine-tuning (Online RFT)}: Unlike standard RFT, Online RFT commences with the SFT-initialized policy, but iteratively updates the model via fine-tuning on outputs generated in real-time by the current policy.
    \item \textbf{PPO/GRPO}: PPO/GRPO also begins with the SFT model and strengthens it by reinforcing on-the-fly outputs generated by the current policy model.
\end{itemize}

The elements comprising these approaches are outlined in Table \ref{tab:data_policy}. For an in-depth explanation of the derivation procedures, please consult Appendix \ref{app:analysis-rl}.

\paragraph{Analysis of Data Source} The data sources utilized are categorized into two primary types: online sampling and offline sampling. Online sampling refers to gathering training data through the exploratory behaviors of the policy model being updated in real-time, whereas offline sampling entails using data generated by the initial SFT model. Specifically, RFT and DPO are representative of the offline sampling paradigm, while Online RFT and GRPO adopt the online sampling approach.

Figure \ref{fig:rl-analysis} illustrates that Online RFT notably exceeds the performance of RFT across both evaluated benchmarks. At the outset of training, the effectiveness of Online RFT is nearly indistinguishable from that of RFT; however, as training advances, Online RFT establishes a pronounced performance lead, showcasing the benefits of online training. This observation aligns with expectations, as the actor and the SFT model initially produce highly similar outputs, and the data differences remain subtle. As training proceeds, the actor samples become more distinct, making the use of data collected directly from the evolving policy increasingly advantageous.

\paragraph{Analysis of Gradient Coefficient} The transformation of reward signals into gradient coefficients is a key step for parameter updates in these algorithms. In our experimental setup, the reward function is specified as either `Rule' or `Model'. The `Rule' setting evaluates the appropriateness of responses through answer correctness, while `Model' involves training a reward model to assign scores to responses; this model is trained on data evaluated by rules. As detailed in Equations \ref{eq:GC-RFT} and \ref{eq:GC-GRPO}, GRPO and Online RFT differ in a critical aspect: GRPO modulates its gradient coefficient dynamically according to the magnitude of the reward assigned by the reward model, thereby providing tailored reinforcement or penalties. In contrast, Online RFT does not incorporate such adaptability—it reinforces all correctly answered responses equally and refrains from penalizing incorrect responses.

As illustrated in Figure \ref{fig:rl-analysis}, GRPO outperforms online RFT, which emphasizes the effectiveness of modifying positive and negative gradient coefficients. Moreover, GRPO+PS achieves better results than GRPO+OS, underscoring the advantages provided by step-aware, fine-grained gradient coefficients. Additionally, we examine the impact of iterative RL, conducting two cycles of iteration during our experiments. According to Figure \ref{fig:iter-rl}, iterative RL produces substantial gains, with the most notable improvement arising after the first iteration.

\subsubsection{Why RL Works?}  
This study utilizes reinforcement learning on a selected subset of instruction tuning data and observes that it yields marked improvements over the instruction-tuned baseline. To elucidate the mechanism by which reinforcement learning contributes to these gains, we compare Pass@K and Maj@K accuracies of the Instruct and RL-enhanced models on two evaluation benchmarks. Figure \ref{fig:maj-pass} reveals that while RL increases Maj@K scores, Pass@K remains unchanged. This suggests that the advantage brought by RL primarily lies in fortifying the robustness of the output distribution, i.e., \textbf{the boost in performance stems from raising the ranking of correct answers among the TopK, not from augmenting the underlying abilities of the model.} In line with this, \citep{wang2023making} identify a \textbf{misalignment problem} associated with reasoning tasks in SFT models, showing that applying preference alignment techniques \citep{yuan2023rrhf,song2023preference,wang2023making} can enhance SFT model reasoning performance.

\subsubsection{How to Achieve More Effective RL?}
\label{sec:effec_rl}
Our experiments confirm that RL methods perform particularly well on mathematical reasoning benchmarks. We also propose a unified framework that accommodates various prominent training paradigms, viewing them all as manifestations or simplifications of RL. As described in Equation \ref{eq:objective}, this framework is characterized by three primary factors: Data Source, Algorithm, and Reward Function. We discuss several prospective avenues for advancing each of these elements.

\paragraph{Data Source} The data source forms the foundation for all training approaches. Within the context of RL, this term denotes unlabeled queries, accompanied by responses generated through sampling from the policy model. For this work, we exclusively sample questions from the instruction tuning dataset and use basic nucleus sampling to generate outputs, which may contribute to RL improvements being confined to Maj@K. Looking forward, we aim to extend the RL pipeline to include out-of-distribution prompts, along with \textbf{advanced sampling (decoding) strategies} such as tree-search-based methods \citep{yao2023tree}. Furthermore, \textbf{efficient inference techniques} \citep{xia-etal-2023-speculative,leviathan2023fast,kwon2023efficient,xia2024unlocking}—critical for the efficient exploration by policy models—will be increasingly significant.

\paragraph{Algorithms} Algorithms utilize the incoming data and reward feedback to adjust the gradient coefficients for updating model parameters. Referring to Equation \ref{eq:objective}, contemporary methods generally place significant \textbf{TRUST} in the reward function’s feedback, using it to raise or lower the conditional probability assigned to specific tokens. Nevertheless, the reliability of the reward signal cannot always be taken for granted, particularly in tasks characterized by high complexity. For instance, even within the PRM800K datasets \citep{lightman2023let}, which feature detailed annotation from skilled annotators, there remains an estimated error rate of 20\% in the annotations\footnote{\url{https://github.com/openai/prm800k/issues/12\#issuecomment-1728491852}}. Given these limitations, we aim to investigate reinforcement learning algorithms that can maintain robustness when faced with noisy reward information. Our perspective is that alignment strategies inspired by the \textbf{WEAK-TO-STRONG} framework \citep{burns2023weak} have the potential to fundamentally transform learning algorithms.

\paragraph{Reward Function} The reward function serves as the primary origin of the learning signal during training. In reinforcement learning settings, this function is typically implemented as a neural reward model. We propose that reward models should address three pivotal aspects: 1) \textbf{Enhancing the generalization capabilities of the reward model.} The reward model needs to generalize effectively, managing both out-of-distribution queries and sophisticated decoding results; otherwise, reinforcement learning may merely reinforce the existing distribution of LLMs without elevating their core abilities; 2) \textbf{Capturing and representing the uncertainty inherent in the reward model.} Accurately quantifying uncertainty may enable a meaningful interaction between imperfect reward models and weak-to-strong learning approaches; 3) \textbf{Developing efficient, high-quality process reward models} that are capable of supplying detailed and actionable training feedback throughout the reasoning process \citep{lightman2023let,wang2023math}.

\section{Conclusion, Limitation, and Future Work} 

In this work, we introduce DeepSeekMath, which surpasses all available open-source models on the challenging MATH benchmark and achieves performance comparable to that of proprietary systems. DeepSeekMath~begins with DeepSeek-Coder-v1.5 7B as its foundation and is trained continually on 500B tokens, including an essential 120B math-specific tokens largely extracted from Common Crawl. Our thorough ablation experiments demonstrate that web-sourced data constitute a valuable resource for curating high-quality mathematical corpora, whereas the contributions from arXiv are less impactful than anticipated. We propose Group Relative Policy Optimization (GRPO), an adaptation of Proximal Policy Optimization (PPO), which boosts mathematical reasoning performance with a reduction in memory requirements. Our empirical findings confirm that GRPO is advantageous, maintaining effectiveness even when DeepSeekMath-Instruct 7B has already attained high benchmark scores. Furthermore, we establish a unified framework for interpreting diverse methods and highlight several avenues for advancing reinforcement learning techniques.

Despite its strong performance on quantitative reasoning datasets, DeepSeekMath~exhibits comparatively lower proficiency in areas such as geometry and theorem-proving, trailing behind some closed models. As observed in preliminary tests, the model struggles with triangle- and ellipse-related questions, hinting at possible biases in training and fine-tuning data selection. Moreover, due to limitations in model size, DeepSeekMath~lags behind GPT-4 with respect to few-shot learning abilities: while GPT-4 benefits from few-shot prompts, DeepSeekMath~demonstrates minimal improvement over zero-shot performance under similar conditions. Looking ahead, our focus will be on further refining our data selection methodology to enhance the quality of pre-training datasets. Additionally, we plan to investigate the directions described in Section \ref{sec:effec_rl} for developing more effective reinforcement learning approaches for large language models.

\end{document}